\section{Discussion}\label{sec:discussion}
% Budget: 1 page together with the conclusion. Draft; to be revised once the evaluation results are in.

\paragraph{Implications for Architects.} The results turn the placement of monitors from a design heuristic into a consequence of the contract hierarchy: once the contracts of the components refine the contract of the system, one guarantee monitor per component and one assumption monitor at the inputs of the system suffice, and each monitor needs access to the interface of its component only. This suits components that are supplied by third parties or learned from data, whose internals are not accessible. Extendability adds an obligation that is easily overlooked: a refinement proof is of no use at runtime if it rests on guarantees that cannot be upheld forever. Theorem~\ref{thm:deadend} makes this obligation checkable on finite prefixes.

\paragraph{Limitations.} The contract language is propositional and restricted to the past and to bounded windows. Properties of individual objects over time, aggregates other than counting, and liveness are outside its scope (Sect.~\ref{sec:synthesis}). The guarantees a learned component can give are statistical; they are not proved but stated as contracts, and the theory tells what follows if they hold and who is to blame if they do not. Whether a contract is adequate, \ie, whether its violation indicates a hazard and its satisfaction the absence of one, remains a matter of validation. Finally, the monitors are executed on the runtime of Lean, which is not designed for hard real-time constraints.

\paragraph{Threats to Validity.} 1)~\emph{Internal Validity:} The mechanised theorems hold for the formal model. The adapter that samples the topics in each interval, the fixed duration of an interval, lost or reordered messages, and floating-point arithmetic in atoms are not part of the model. \todo{State how the adapter was validated.} 2)~\emph{External Validity:} The evaluation uses one perception pipeline and recorded data in open loop; closed-loop effects and other architectures are not covered. 3)~\emph{Construct Validity:} Faults injected at the level of topics approximate, but do not reproduce, the failures of learned components. The rating of requirements for RQ1 involves judgement; \todo{report who rated and how disagreements were resolved.} 4)~\emph{Conclusion Validity:} The requirements of the design set shaped the contract language. We counter this bias by reporting a held-out set separately (Sect.~\ref{sec:evaluation}).
